\chapter{Introduction}
\label{ch:introduction}

\section{Background and Context}

Skin cancer represents one of the most common malignancies worldwide, with melanoma being the deadliest form. Early detection and accurate classification are critical factors determining treatment outcomes and patient survival rates. Dermatologists traditionally rely on visual examination and dermoscopic analysis, a practice that requires extensive training and experiences significant inter-observer variability \cite{esteva2019dermatologist}.

The emergence of digital pathology and computational imaging has opened new avenues for automated skin lesion analysis. Deep learning models, trained on large annotated datasets of dermatoscopic images, have demonstrated remarkable potential in assisting clinicians with diagnostic decisions. However, achieving accurate, reliable, and clinically applicable automated diagnosis systems remains a significant challenge due to the high visual similarity among different skin lesion types and the subtle morphological differences that distinguish benign lesions from malignant melanomas.

\section{Motivation and Significance}

The motivation for this research stems from several critical observations:

\begin{enumerate}
    \item \textbf{Diagnostic Accuracy Gap}: Studies have shown that even experienced dermatologists exhibit diagnostic accuracy of 75-95\% for skin lesion classification, with significant variation depending on expertise and lesion characteristics.
    
    \item \textbf{Scalability Challenges}: The global shortage of qualified dermatologists makes immediate access to expert diagnosis impossible for millions of patients in developing regions.
    
    \item \textbf{Computational Advances}: Recent breakthroughs in vision transformer architectures, particularly the Swin Transformer, have demonstrated superior performance characteristics compared to traditional convolutional neural networks.
    
    \item \textbf{Clinical Applicability}: There is urgent need for computer-aided diagnosis (CAD) systems that can integrate seamlessly into clinical workflows and provide reliable decision support.
\end{enumerate}

\section{Research Objectives}

The primary objectives of this thesis are:

\begin{itemize}
    \item To design and implement a state-of-the-art Swin Transformer-based deep learning model optimized specifically for multi-class skin lesion classification on the HAM10000 dataset.
    
    \item To conduct comprehensive experimental evaluation comparing the proposed model against baseline approaches including traditional CNNs and other transformer variants.
    
    \item To perform ablation studies and sensitivity analyses to understand the contribution of different architectural components and training strategies.
    
    \item To demonstrate clinical feasibility through rigorous validation metrics aligned with medical diagnostic requirements.
    
    \item To provide actionable insights for practitioners and researchers regarding deployment considerations and practical implementation strategies.
\end{itemize}

\section{Thesis Organization}

This thesis is organized as follows:

\begin{itemize}
    \item \textbf{Chapter 2} presents a comprehensive literature review of skin cancer diagnosis methods, deep learning architectures, and transformer-based models.
    
    \item \textbf{Chapter 3} formally defines the problem statement and articulates the research motivation and gaps.
    
    \item \textbf{Chapter 4} describes the HAM10000 dataset, data preprocessing, and augmentation strategies.
    
    \item \textbf{Chapter 5} presents the proposed Swin Transformer architecture and methodology.
    
    \item \textbf{Chapter 6} details the implementation framework, hyperparameter configuration, and training procedures.
    
    \item \textbf{Chapter 7} describes the experimental design and validation approaches.
    
    \item \textbf{Chapter 8} presents comprehensive experimental results with detailed analysis.
    
    \item \textbf{Chapter 9} provides in-depth discussion of findings, implications, and limitations.
    
    \item \textbf{Chapter 10} concludes the thesis and outlines directions for future research.
\end{itemize}

\section{Contributions}

The key contributions of this research include:

\begin{enumerate}
    \item A novel application of Swin Transformer architecture to skin cancer diagnosis with custom optimization for medical imaging constraints.
    
    \item Comprehensive empirical evaluation demonstrating superior performance compared to existing approaches.
    
    \item Detailed ablation studies providing insights into architectural design principles for medical imaging tasks.
    
    \item Practical implementation guidelines and considerations for clinical deployment.
    
    \item Open-source implementation framework facilitating future research and clinical adoption.
\end{enumerate}
